\documentclass[a4paper,oneside,12pt]{extreport}

\usepackage{mmap}
\usepackage[T2A]{fontenc}
\usepackage[utf8]{inputenc}
\usepackage[english,russian]{babel}

% Текст отчёта следует печатать, соблюдая следующие размеры полей:
% левое — 30 мм, правое — 15 мм, верхнее и нижнее — 20 мм.
\usepackage[left=20mm, right=15mm, top=15mm, bottom=15mm]{geometry}

% \setlength{\parindent}{1.25cm} % Абзацный отступ

\usepackage{setspace}
%\onehalfspacing % Полуторный интервал

\frenchspacing % Равномерные пробелы
\usepackage{indentfirst} % Красная строка

\usepackage{microtype}
\sloppy

\usepackage{titlesec}
\titlespacing*{\chapter}{0pt}{-30pt}{8pt}
\titlespacing*{\section}{\parindent}{*4}{*4}
\titlespacing*{\subsection}{\parindent}{*4}{*4}
\titleformat{\chapter}{\LARGE\bfseries}{\thechapter}{20pt}{\LARGE\bfseries}
\titleformat{\section}{\Large\bfseries}{\thesection}{40pt}{\Large\bfseries}

\usepackage{graphicx}
\usepackage{caption}

\usepackage[unicode,pdftex]{hyperref}
\hypersetup{hidelinks}

%% title begin
\usepackage{wrapfig}

\makeatletter
	\def\vhrulefill#1{\leavevmode\leaders\hrule\@height#1\hfill \kern\z@}
\makeatother
%% title end

%% begin code
\usepackage{listings}
\usepackage{xcolor}

\lstset{
	basicstyle=\footnotesize\ttfamily,
	breakatwhitespace=true,
	breaklines=true,
	commentstyle=\color{gray},
	frame=single,
	keywordstyle=\color{blue},
	showstringspaces=false,
	stringstyle=\color{red},
	tabsize=8
}

\lstdefinestyle{lispinline}{
	frame=none,
	language=Lisp
}

\newcommand{\code}[1]{\texttt{#1}}
%% end code

%% begin theorem
\usepackage{amsthm}

\makeatletter
\newtheoremstyle{indented}
	{}% measure of space to leave above the theorem
	{}% measure of space to leave below the theorem
	{}% name of font to use in the body of the theorem
	{\parindent}% measure of space to indent
	{\bfseries}% name of head font
	{.}% punctuation between head and body
	{ }% space after theorem head; " " = normal interword space
	{}% header specification (empty for default)
\makeatother

\theoremstyle{indented}

\newtheorem{definition}{Определение}[section]
\newtheorem{example}{Пример}[section]
\newtheorem{theorem}{Теорема}[section]
\newtheorem{task}{Задание}

\makeatletter
\DeclareRobustCommand\bfseriesitshape{%
	\not@math@alphabet\itshapebfseries\relax
	\fontseries\bfdefault
	\fontshape\itdefault
	\selectfont
}
\makeatother

\DeclareTextFontCommand{\textbfit}{\bfseriesitshape}
\DeclareTextFontCommand{\define}{\bfseriesitshape}
%% end theorem

%% begin columns
\usepackage{etoolbox,refcount}
\usepackage{multicol}

\newcounter{countitems}
\newcounter{nextitemizecount}
\newcommand{\setupcountitems}{%
	\stepcounter{nextitemizecount}%
	\setcounter{countitems}{0}%
	\preto\item{\stepcounter{countitems}}%
}
\makeatletter
\newcommand{\computecountitems}{%
	\edef\@currentlabel{\number\c@countitems}%
	\label{countitems@\number\numexpr\value{nextitemizecount}-1\relax}%
}
\newcommand{\nextitemizecount}{%
	\getrefnumber{countitems@\number\c@nextitemizecount}%
}
\newcommand{\previtemizecount}{%
	\getrefnumber{countitems@\number\numexpr\value{nextitemizecount}-1\relax}%
}
\makeatother
\newenvironment{AutoMultiColItemize}{%
	\ifnumcomp{\nextitemizecount}{>}{3}{\begin{multicols}{2}}{}%
		\setupcountitems\begin{itemize}}%
		{\end{itemize}%
		\unskip\computecountitems\ifnumcomp{\previtemizecount}{>}{3}{\end{multicols}}{}}
\makeatother
\newenvironment{AutoMultiColEnumerate}{%
	\ifnumcomp{\nextitemizecount}{>}{3}{\begin{multicols}{2}}{}%
		\setupcountitems\begin{enumerate}}%
		{\end{enumerate}%
		\unskip\computecountitems\ifnumcomp{\previtemizecount}{>}{3}{\end{multicols}}{}}
%% end columns

\usepackage{array}
\usepackage{float}
\restylefloat{table}

\usepackage{amsmath}
\usepackage{amssymb}

\begin{document}

\begin{titlepage}
	{\large % 14pt instead of 12pt
	\onehalfspacing
	\centering

	\begin{wrapfigure}[7]{l}{0.14\linewidth}
		\vspace{3mm}
		\hspace{-10mm}
		\includegraphics[width=0.93\linewidth]{inc/img/bmstu-logo}
	\end{wrapfigure}
	{\singlespacing \footnotesize \bfseries Министерство науки и высшего образования Российской Федерации\\Федеральное государственное бюджетное образовательное учреждение\\высшего образования\\<<Московский государственный технический университет\\имени Н.~Э.~Баумана\\ (национальный исследовательский университет)>>\\(МГТУ им. Н.~Э.~Баумана)\\}

	\vspace{-2.2mm}
	\vhrulefill{0.9mm}\\
	\vspace{-7.5mm}
	\vhrulefill{0.2mm}\\
	\vspace{2mm}

	{\doublespacing \small \raggedright ФАКУЛЬТЕТ \hspace{25mm} «Информатика и системы управления»\\
	КАФЕДРА \hspace{5mm} «Программное обеспечение ЭВМ и информационные технологии»\\}

	\vspace{30mm}

	\textbf{ОТЧЁТ}\\
	По лабораторной работе №6\\
	По курсу: «Функциональное и логическое программирование»\\
	%Тема: «Тема работы»\\

	\vspace{60mm}

	\hspace{70mm} Студент:       \hfill Керимов~А.~Ш.\\
	\hspace{70mm} Группа:        \hfill ИУ7-64Б\\
	\hspace{70mm} Преподаватели: \hfill Толпинская~Н.~Б.,\\
	                             \hfill Строганов~Ю.~В.\\

	\vfill

	Москва\\
	\the\year\\}
\end{titlepage}

\setcounter{page}{2}


\begin{task}
	Написать функцию, которая по своему списку-аргументу \code{lst} определяет, является ли он палиндромом (то есть равны ли \code{lst} и \code{(reverse lst)}).

	\begin{lstlisting}[language=Lisp, gobble=16]
		(defun is_palindrom (lst)
		  (and (equal (null lst) nil)
		       (equal (reverse lst) lst)))
	\end{lstlisting}
\end{task}

\begin{task}
	Написать предикат \code{set-equal}, который возвращает \code{t}, если два его множества-аргумента содержат одни и те же элементы, порядок которых не имеет значения.
	\begin{lstlisting}[language=Lisp, gobble=16]
		(defun set-equal1 (a b)
		  (and (null (set-difference a b))
		       (null (set-difference b a))))

		(defun set-equal2 (a b)
		  (not (set-exclusive-or a b)))
	\end{lstlisting}
\end{task}

\begin{task}
	Напишите функцию \code{swap-first-last}, которая переставляет в списке-аргументе первый и последний элементы.

	\begin{lstlisting}[language=Lisp, gobble=16]
		(defun swap-first-last (x)
		  (cond ((null (cdr x)) x)
		        (t (append (last x) (cdr (butlast x)) (list (car x))))))
	\end{lstlisting}
\end{task}

\begin{task}
	Напишите функцию \code{swap-two-elements}, которая переставляет в списке- аргументе два указанных своими порядковыми номерами элемента в этом списке.

	\begin{lstlisting}[language=Lisp, gobble=16]
		(defun swap-two-elements1 (x i1 i2)
		  (cond ((null (rotatef (nth i1 x) (nth i2 x))) x)))

		(defun swap-two-elements2 (x i1 i2)
		  (let ((i (min i1 i2)) (j (max i1 i2)))
		       (append
		         (subseq x 0 i)
		         (list (nth j x))
		         (subseq x (+ i 1) j)
		         (list (nth i x))
		         (subseq x (+ j 1)))))

		(defun swap-two-elements3 (x i1 i2)
		  (let ((i (min i1 i2)) (j (+ (max i1 i2) 1)))
		    (append
		      (subseq x 0 i)
		      (swap-first-last (subseq x i j))
		      (subseq x j))))
	\end{lstlisting}
\end{task}

\begin{task}
	Напишите две функции, \code{swap-to-left} и \code{swap-to-right}, которые производят круговую перестановку в списке-аргументе влево и вправо, соответственно.

	\begin{lstlisting}[language=Lisp, gobble=16]
		(defun swap-to-left (lst &optional (k 1))
		  (let ((k (rem k (length lst))))
		       (append (subseq lst k) (subseq lst 0 k))))

		(defun swap-to-right (lst &optional (k 1))
		  (setf len (length lst))
		  (let ((k (- len (rem k len))))
		       (append (subseq lst k) (subseq lst 0 k))))
	\end{lstlisting}
\end{task}

\begin{task}
	Напишите функцию, которая умножает на заданное число-аргумент все числа из заданного списка-аргумента, когда
	a) все элементы списка — числа,
	б) элементы списка — любые объекты.

	\begin{lstlisting}[language=Lisp, gobble=16]
		(defun mul_n (lst n)
		  (mapcar (lambda (x) (* n x)) lst))

		(defun mul_all (lst n)
		  (mapcar
		    (lambda (x)
		      (cond
		        ((numberp x) (* n x))
		        ((listp x) (mul_all x n))
		        (t x)))
		    lst))
	\end{lstlisting}
\end{task}

\begin{task}
	Напишите функцию, \code{select-between}, которая из списка-аргумента,
	содержащего только числа, выбирает только те, которые расположены между двумя указанными границами-аргументами и возвращает их в виде списка (упорядоченного по возрастанию списка чисел).

	\begin{lstlisting}[language=Lisp, gobble=16]
		(defun select-between (lst a b)
		  (let ((a (min a b)) (b (max a b)))
		       (sort
		         (remove-if-not (lambda (x)
		           (and (<= a x) (<= x b))) lst)
		         #'<)))
	\end{lstlisting}
\end{task}

\end{document}
